\documentclass{article}
\usepackage{amsmath,amssymb}
\usepackage{xurl}
\usepackage[hidelinks]{hyperref}
% Shared commands for notes in docs/paper-gaps/.

\DeclareUnicodeCharacter{2080}{\ifmmode{}_{0}\else\textsubscript{0}\fi}
\DeclareUnicodeCharacter{2081}{\ifmmode{}_{1}\else\textsubscript{1}\fi}
\DeclareUnicodeCharacter{2082}{\ifmmode{}_{2}\else\textsubscript{2}\fi}
\DeclareUnicodeCharacter{2099}{\ifmmode{}_{n}\else\textsubscript{n}\fi}

\newcommand{\Lean}{\textsc{Lean}}
\ExplSyntaxOn
\NewDocumentCommand{\leanid}{m}{
  \ifmmode
    \textsf{\detokenize{#1}}
  \else
    \group_begin:
    \urlstyle{sf}
    \tl_set:Nn \l_tmpa_tl {#1}
    \tl_replace_all:Nnn \l_tmpa_tl {\_} {_}
    \exp_args:NV \path \l_tmpa_tl
    \group_end:
  \fi
}
\ExplSyntaxOff
\providecommand{\tr}{\operatorname{tr}}
\newcommand{\tnleanIssueBase}{https://github.com/LionSR/TNLean/issues/}
\newcommand{\tnleanPrBase}{https://github.com/LionSR/TNLean/pull/}
\newcommand{\tnleanDocBase}{https://sirui-lu.com/TNLean/docs/}
\newcommand{\ghissue}[1]{\href{\tnleanIssueBase#1}{GitHub issue \##1}}
\newcommand{\ghpr}[1]{\href{\tnleanPrBase#1}{PR \##1}}
\newcommand{\doclink}[1]{\href{\tnleanDocBase#1.html}{\path{#1}}}

% Dirac notation and the identity operator, as in blueprint/src/macros/common.tex.
% \providecommand keeps note-local definitions authoritative.
\usepackage{dsfont}
\providecommand{\ket}[1]{|#1\rangle}
\providecommand{\bra}[1]{\langle #1|}
\providecommand{\braket}[2]{\langle #1 | #2 \rangle}
\providecommand{\ketbra}[2]{|#1\rangle\!\langle #2|}
\providecommand{\Id}{\mathds{1}}
\providecommand{\one}{\mathds{1}}
\providecommand{\C}{\mathbb{C}}
\providecommand{\MN}[1]{M_{#1}(\C)}
\providecommand{\MD}{M_D(\C)}

% Verdict marker: \gapnote{<kind>}{<status>}. Kind is the policy
% classification (clarification, local-correction, scope-restriction,
% unfaithful, false-source, open-gap); status is open, wip (elimination
% underway), resolved, or historical. Machine-read by the site index;
% prints a verdict line.
\newcommand{\gapnote}[2]{\par\noindent\textbf{Verdict:} #1 (#2).\par}

\title{Long-Range Matrix Product States with Measurements: Scope of the Formal Statement}
\author{}
\date{2026-10-02}
\begin{document}
\maketitle
\gapnote{scope-restriction}{open}

\section{What the source states}
The paragraph ``Long-range MPS using measurements'' of~\cite{Malz2024Preparation}
writes the fixed point of a translation-invariant MPS that is not normal, up to
an isometry, as
\begin{align}
    \ket{\Omega'}
     & =\sum_{j=1}^b\alpha_j^{(N)}\bigotimes_{i=1}^{N/q}\ket{\omega_j}_{R_iL_{i+1}}
    =W^{\otimes N/q}\ket{\chi_{N/q}},\qquad
    \ket{\chi_M}=\sum_{j=1}^b\alpha_j\ket j^{\otimes M},
    \label{eq:mswc24_meas_fixed_point}
\end{align}
with $W\colon\ket j\mapsto\ket{\omega_j}$, and prepares it as follows: ``First
create $\ket{\chi_{N/q}}$, which can be done in constant depth with
measurements \ldots\ Subsequently, apply in parallel the isometries $W$ \ldots,
which also takes constant depth.'' With the approximation of the target by a
state of the form (\ref{eq:mswc24_meas_fixed_point}) up to local isometries,
obtained by blocking $q\propto\log(N/\varepsilon)$ sites, it concludes that if
``measurements are only used for the preparation of $\ket{\chi_{N/q}}$, the
depth is $O(\log(N/\varepsilon))$'', for all translation-invariant MPS, short-
or long-range correlated.

\section{The statements formalized}
\paragraph{Orthogonal blocks and block lengths dividing $N$.}
Let $A^i=\bigoplus_j\operatorname{diag}(\mu_{j,1},\dots,\mu_{j,m_j})\otimes A_j^i$
with normal blocks $A_j$ in the gauge of eq.~(5) of the source and nonzero
weights. The first formal statement prepares, with measurements, a unit vector
of error at most $\varepsilon$ in depth at most $c\log(N/\varepsilon)$, for
$N\geq2$, $0<\varepsilon\leq1$, and block lengths $q$ with
$a\log(N/\varepsilon)+b'\leq q\leq2(a\log(N/\varepsilon)+b')$ under three
restrictions.\footnote{Formal declarations:
    \leanid{MPSPreparation.exists_isPreparedWithMeasurementsAndCircuitInDepth_le_log_repeatedBlockSum}
    and
    \leanid{MPSPreparation.exists_isPreparedWithMeasurementsAndCircuitInDepth_copyApproxVector}
    in \doclink{TNLean/MPS/Preparation/NonNormalMeasurementPreparation}.}
\begin{enumerate}
    \item The $q$-site states of distinct blocks are orthogonal. Without this the
          block form of the positive part fails
          \cite{gap:mswc24_block_form_mixed_overlap}, and the isometry $V$ of the
          blocked direct sum is no longer the sum of the isometries $V_j$ of the
          blocks, which is what the block unitary of this statement implements.
    \item The block length $q$ divides $N$, so that the chain consists of
          $N/q$ blocks of $q$ sites, as in eq.~(10) of the source. Together with
          the window above, this statement applies only to chain lengths $N$
          with a divisor in this window, and it is vacuous for the others, for
          instance for every large prime $N$.
    \item The state prepared is the corrected approximating state of
          \cite{gap:mswc24_repeated_block_corrected_state}, which agrees with
          the state of eq.~(S7) when every multiplicity is one and every weight
          is positive.
\end{enumerate}

\paragraph{Overlapping blocks and every chain length.}
The second formal statement removes the three restrictions above.\footnote{Formal
    declarations:
    \leanid{MPSPreparation.exists_isPreparedWithMeasurementsAndCircuitInDepth_le_log_of_mpvState_ne_zero}
    and
    \leanid{MPSPreparation.exists_isPreparedWithMeasurementsAndCircuitInDepth_le_log}
    in \doclink{TNLean/MPS/Preparation/OverlappingMeasurementPreparation}.}
For $A$ as above, with every fixed point positive definite and the mixed
transfer maps of distinct blocks of spectral radius below one, there is $c$
such that for every $N\geq2$ and $0<\varepsilon\leq1$ at which
$\ket{\phi_N(A)}\neq0$, a unit vector of error at most $\varepsilon$ is prepared
with measurements and a circuit in depth at most $c\log(N/\varepsilon)$. There
is $N_0$ such that $\ket{\phi_N(A)}\neq0$ for every $N\geq N_0$ at which the
weights $\beta_j=\sum_k\mu_{j,k}^N$ are not all zero. The $q$-site states of
distinct blocks may overlap, the multiplicities and the complex weights are
arbitrary, and the chain is cut into blocks ``all of the same size, $q_N$,
except for the last one, which may be larger'', as in the proof of Theorem~1 of
the source. The statement reads the construction as follows.
\begin{enumerate}
    \item The isometries applied on the blocks are the isometric factors
          $V^{(\ell)}$ of the polar decompositions of the blocked maps of the
          direct sum $\bigoplus_jA_j$ with weights one and one copy of each
          block, not that of the blocked map of $A$, whose positive part does
          not have the block form (S5) when the blocks overlap or repeat. For
          large $\ell$ this blocked map is injective on the bond pairs of the
          blocks, by the positive definiteness of the limit of its Gram matrix,
          so $V^{(\ell)}$ is an isometry on these pairs. The state prepared is
          \begin{align*}
              \ket\psi=\sum_j\alpha_j\Bigl(\bigotimes_kV^{(\ell_k)}\Bigr)
              \bigotimes_k\ket{\omega_j},\qquad
              \alpha_j=\frac{\beta_j}{(\sum_l\lvert\beta_l\rvert^2)^{1/2}},
          \end{align*}
          with $\omega_j$ the pair of the fixed point of block $j$. The weights
          enter only through the amplitudes, so the overlap
          $\braket\psi{\phi_N(A)}$ is a quadratic form
          $\sum_{j,j'}\overline{\beta_j}\beta_{j'}z_{jj'}$ divided by
          $(\sum_l\lvert\beta_l\rvert^2)^{1/2}$, with
          $\lvert z_{jj'}-\delta_{jj'}\rvert\leq Cy\,e^{Cy}$ and
          $y=Me^{-\gamma q/\xi}$. The error is therefore at most
          $CMe^{-\gamma q/\xi}$ with $C$ independent of the weights, and the
          factor $\min(1,\sum_j\lvert\beta_j\rvert^2)^{-1/2}$ of the error bound
          of the corrected state \cite{gap:mswc24_repeated_overlap_small_weight_sum}
          does not enter.\footnote{Formal declaration:
              \leanid{MPSTensor.exists_one_sub_norm_inner_sum_blockIsometryState_le}
              in \doclink{TNLean/MPS/Preparation/OverlappingBlockStates}.}
    \item The rate is $e^{-\gamma q/\xi}$ with $\xi$ bounding the correlation
          lengths of the blocks and of the mixed transfer maps of distinct
          blocks \cite{gap:mswc24_block_form_mixed_overlap}.
    \item The periodic state is assumed nonzero, as the normalization of the
          source presupposes \cite{gap:mswc24_depth_upper_bound_nonzero_state};
          this holds for $N\geq N_0$ when the weights $\beta_j$ are not all
          zero.
\end{enumerate}
For $N$ below the block length the normalized target itself is prepared, as
the state of one block or, below a fixed length, by a circuit of bounded depth.

\section{Two readings of the construction}
\paragraph{The circuit after the measurement.}
The model of preparation with measurements of~\cite{Piroli2021LOCC} applies a
circuit, then measurements and local corrections. The construction of the
source applies circuits after the preparation of $\ket{\chi_{N/q}}$. The formal
statement counts such a preparation as a vector prepared with measurements in
depth $T_1$, followed by a local circuit of depth $T_2$ that does not depend on
the outcomes, of total depth $T_1+T_2$.

\paragraph{The depth of the GHZ-type state.}
The state $\ket{\chi_{N/q}}$ lives on the $N/q$ blocks. On the chain of $N$
sites, the label of block $k$ is stored in the last $r_1$ sites of the block,
and the first $r_1$ sites of the next block serve as the ancilla of Example~1
of~\cite{Piroli2021LOCC}; this ancilla is $q-r_1$ sites away from the register
of its own block, and the controlled shift between them is carried out with
SWAP gates in depth $O(q)$. The formal statement therefore prepares
$\ket{\chi_{N/q}}$ in depth $O(q)$ rather than in constant depth, a scope
restriction of this step. Constant depth
on the chain of $N$ sites would need the controlled shifts between distant
sites to be implemented with further measurements, for instance by
teleportation along pairs of neighbouring sites. Since the isometries of the
blocked tensor already take depth $O(q)$, the total depth $O(\log(N/\varepsilon))$
is unaffected.

\section{Open}
The first two restrictions of the first statement are removed by the second
statement, and the third does not apply to it. What remains open is the depth
of the GHZ-type state: both statements prepare $\ket{\chi_{N/q}}$ in depth
$O(q)$ on the chain of $N$ sites, as described above, rather than in constant
depth. This does not change the total depth $O(\log(N/\varepsilon))$. Constant
depth needs the controlled shifts between the register of a block and the
ancilla at the start of the next block to be carried out with further
measurements, for instance by teleportation along pairs of neighbouring sites.

\bibliographystyle{alpha}
\bibliography{references}
\end{document}
